\chapter{Symbol Resolution, Archives, and Separate Compilation}
\label{ch:link-symbols}

Two object files can contain perfectly valid machine instructions and still fail to form a program.
One may reference a function that neither defines.
Both may define the same global name with incompatible intentions.
A library may contain the missing definition but appear at a point in the command where the linker's extraction policy does not select it.
These are not address-arithmetic failures; they arise before layout from the meaning of the set of objects that participates in the link.

This chapter develops that meaning explicitly.
Chapter~\ref{ch:linking-loading} provides the introductory pipeline; the following four chapters expand its stages and culminate in a working Lua-mapped static linker.
The ELF symbol specification defines local, global, weak, and special-section states; production linker documentation defines additional search and policy choices~\cite{xinuos-2026-elf,gnu-2026-ld}.

\section{Local spelling is not global identity}
\label{sec:ch57-identity}

An assembler may emit a private label named \code{loop} in each object.
The linker must not merge those two labels merely because their spellings match.
A local symbol's identity includes its owning object and symbol index.
A relocation to that local symbol resolves within that object, even when another object's local symbol has the same name and offset.
Figure~\ref{fig:ch57-local} makes the ownership boundary visible.

\begin{figure}[htbp]
\centering
\begin{tikzpicture}[node distance=17mm and 22mm]
\node[basic block] (a) {alpha.o relocation};
\node[basic block,right=of a] (b) {beta.o relocation};
\node[basic block,below=of a] (al) {alpha.o local loop};
\node[basic block,below=of b] (bl) {beta.o local loop};
\draw[cfg edge] (a)--(al); \draw[cfg edge] (b)--(bl);
\draw[dashed] (al)--node[below=3mm,font=\small] {same spelling; distinct identities}(bl);
\end{tikzpicture}
\caption{Local references remain object-qualified; a global name table must not absorb them.}
\label{fig:ch57-local}
\end{figure}

Global symbols introduce a different lookup space.
Each participating object can contribute declarations or definitions to the global name's resolution state.
That state is not just the first numeric address encountered.
It records which definitions compete, which references remain unsatisfied, and which binding rules apply.
Final addresses are usually assigned only after this choice has selected an owning section and offset.

The case-study linker uses a direct scan of participating symbol tables to select definitions.
Its local path bypasses global lookup entirely.
This is slower than a production hash table but makes the namespace distinction straightforward to inspect and test.

\section{Resolve binding before assigning addresses}
\label{sec:ch57-binding}

For the ordinary defined-global and defined-weak cases, one strong global definition wins over weak definitions with the same name.
Two strong definitions are an error under the case-study policy.
When only weak definitions compete, the case study chooses the first in command-line object order.
That tie policy is explicit rather than inferred from accidental hash iteration.
Unresolved weak references evaluate to zero, while unresolved strong references fail the static link~\cite{xinuos-2026-elf}.

Table~\ref{tab:ch57-binding} shows the implemented subset.
It deliberately omits common symbols, unique symbols, symbol versions, and dynamic preemption; those mechanisms require additional states and rules.
Trying to fit them into a Boolean \code{defined} flag is a common source of precedence bugs.

\begin{table}[htbp]
\caption{The static case study selects one definition using binding, not the order of the first reference.}
\label{tab:ch57-binding}
\centering
\begin{tabular}{lll}
\toprule
Existing choice & New definition & Result \\
\midrule
None & weak & choose weak \\
None & strong & choose strong \\
Weak & weak & retain first weak \\
Weak & strong & replace with strong \\
Strong & weak & retain strong \\
Strong & strong & duplicate-definition error \\
\bottomrule
\end{tabular}
\end{table}

\begin{algorithm}[htbp]
\caption{Select a definition for one global name in an ordered object set.}
\label{alg:ch57-resolve}
\begin{pseudocode}
\Require all participating objects, one nonlocal name, strong/weak binding subset
\Ensure selected definition or a duplicate-definition diagnostic
\State $chosen\gets none$
\ForAll{defined nonlocal symbols of that name in object order}
  \If{both the new symbol and $chosen$ are strong}
    \State \Return duplicate-definition error
  \EndIf
  \If{$chosen=none$ or a strong symbol replaces a weak choice}
    \State $chosen\gets$ new symbol with its owning object
  \EndIf
\EndFor
\State \Return $chosen$
\end{pseudocode}
\end{algorithm}

Algorithm~\ref{alg:ch57-resolve} is linear in the scanned symbol population for one query.
Calling it independently for every nonlocal symbol and relocation can be quadratic or worse in the combined symbol/reference counts.
A production linker interns global names and accumulates resolution states once, reducing repeated search at the cost of a more involved update mechanism.
The teaching implementation documents its scan-based bound instead of labeling all symbol lookup \enquote{constant time}.

\section{Definitions carry more than binding}
\label{sec:ch57-attributes}

ELF symbol entries also contain type, visibility, size, and section index.
A defined ordinary symbol in a relocatable object holds an offset within its section, not a final process address.
An absolute symbol uses a special section index and keeps its value when sections move.
A common symbol describes storage still to be allocated, with an alignment requirement and size.
An undefined symbol's zero section index does not make it an absolute symbol at address zero.

Common storage illustrates why old language and toolchain conventions remain visible at link time.
Several tentative definitions can request one allocation, often taking the maximum size and strongest alignment under the target linker's policy.
A real initialized definition can supersede that allocation request.
Modern compiler defaults and language rules influence whether the assembler emits common symbols at all.
Our linker rejects common symbols, so it never claims to implement those merging rules by accident.

Visibility is another independent dimension.
A hidden global can participate in resolution within the component being linked while being unavailable for external dynamic lookup.
Protected visibility adds further preemption constraints.
Static resolution and runtime interposition must therefore be described separately.
The case study accepts default and hidden visibility for its fully static output and rejects unsupported visibility encodings.

\section{Archives are demand-selected collections}
\label{sec:ch57-archives}

A static archive is a container of member objects plus indexing information, not one inseparable object.
Linking the archive usually extracts only members selected by unresolved references under the linker's search policy.
Once selected, a member contributes all of its definitions and references, not only the one symbol that caused extraction.
That contribution can create new unresolved names and therefore new extraction opportunities.

This is different from section garbage collection.
Archive extraction chooses which whole objects enter the link.
Section garbage collection later chooses which allocated sections of the selected objects remain live.
If two functions share one selected text section, extraction alone cannot discard one of them.
If they occupy separate sections, a later reachability pass may do so under its root policy.

Use three archive members in two archives.
Archive $A$ contains \code{draw.o}, which defines \code{draw} and references \code{blend}, and \code{helper.o}, which defines \code{helper}.
Archive $B$ contains \code{blend.o}, which defines \code{blend} and references \code{helper}.
The main object initially references \code{draw}.
Table~\ref{tab:ch57-archive} follows a left-to-right search of main, $A$, then $B$.

\begin{table}[htbp]
\caption{A newly selected member can create a dependency on an archive already passed.}
\label{tab:ch57-archive}
\centering
\begin{tabular}{lll}
\toprule
Search event & New selected member & Outstanding names \\
\midrule
Read main & main object & draw \\
Search A & draw.o & blend \\
Search B & blend.o & helper \\
Stop without group rescan & none & helper remains unresolved \\
Rescan A in a group & helper.o & none \\
\bottomrule
\end{tabular}
\end{table}

Figure~\ref{fig:ch57-archive} shows why simply seeing every archive once is insufficient for this order.
The backward dependency is discovered only after $B$ contributes its member.
Some linkers retain more archive information and can provide different behavior; the exact command-line search semantics are tool-specific.
GNU-style archive groups explicitly request repeated searches until extraction reaches closure~\cite{gnu-2026-ld}.

\begin{figure}[htbp]
\centering
\begin{tikzpicture}[node distance=14mm and 19mm]
\node[basic block] (main) {main: needs draw};
\node[basic block,right=of main] (draw) {A: draw.o};
\node[basic block,below=of draw] (blend) {B: blend.o};
\node[basic block,left=of blend] (helper) {A: helper.o};
\draw[cfg edge] (main)--(draw); \draw[cfg edge] (draw)--node[right,font=\small] {needs blend}(blend);
\draw[back edge] (blend)--node[below=3mm,font=\small] {needs helper}(helper);
\end{tikzpicture}
\caption{The archive order and the dependency direction differ, motivating an explicit group-rescan policy.}
\label{fig:ch57-archive}
\end{figure}

\section{Extraction is a monotone fixed point over members}
\label{sec:ch57-closure}

\begin{algorithm}[htbp]
\caption{Search an archive group to member-extraction closure.}
\label{alg:ch57-archive}
\begin{pseudocode}
\Require ordered archives, current definitions and unresolved strong references
\Ensure no unselected member in the group satisfies an outstanding extraction demand
\Repeat
  \State $changed\gets false$
  \ForAll{archives and their members in the declared search order}
    \If{an unselected member defines a demanded name}
      \State mark the member selected before processing its contents
      \State incorporate its definitions using the resolution rules
      \State incorporate its undefined references and recompute demands
      \State $changed\gets true$
    \EndIf
  \EndFor
\Until{no member was selected}
\end{pseudocode}
\end{algorithm}

The monotone set is the set of selected members, not the set of unresolved names.
An unresolved name can disappear when satisfied and a different one can appear from the selected object.
With $M$ members, at most $M$ successful selections occur, so repeated full scans terminate.
If each sweep inspects all member summaries of total size $I$, a coarse bound is $O((M+1)I)$ plus the cost of incorporating selected objects and resolving names.
An archive index and a demand worklist avoid much of the repeated scanning.

Unresolved weak references do not normally trigger archive extraction under the ELF rule~\cite{xinuos-2026-elf}.
That distinction is necessary for optional facilities represented by weak names.
Another subtle consequence is that a weak definition already selected from an ordinary object can satisfy a name before a later archive containing a strong definition is searched.
The strong definition wins only if its member enters the link.
Resolution over selected objects is not equivalent to resolution over every definition physically present in every archive.

\section{Separate compilation preserves a contract, not complete type knowledge}
\label{sec:ch57-abi}

An ELF linker can often confirm that a name is defined without knowing whether callers and callees agree on its language-level type.
The symbol's broad type may say function, but not encode a complete C prototype or calling-convention proof.
Two translation units can agree on spelling while disagreeing on argument classification, structure layout, or return convention.
Successful linking therefore does not establish ABI compatibility.

The demo uses a leaf helper returning an integer in \code{eax} and a startup routine that moves it into \code{edi} for the Linux exit syscall.
It needs no library initialization and does not claim a general C runtime startup sequence.
The function call convention and the system-call convention are different interfaces, even though both use registers.
The AMD64 supplement and operating-system ABI define those details~\cite{matz-2012-amd64}.

Table~\ref{tab:ch57-checks} distinguishes checks available at different layers.
LTO can recover richer IR-level information when all participating toolchains provide compatible representations, but it does not remove the need for object-level ABI and relocation checks.
Assembly objects and prebuilt libraries can remain outside the IR-level view.

\begin{table}[htbp]
\caption{Finding a symbol definition proves less than agreement on the function's complete interface.}
\label{tab:ch57-checks}
\centering
\begin{tabular}{p{0.30\linewidth}p{0.57\linewidth}}
\toprule
Layer & Representative obligation \\
\midrule
Source type checker & Call arguments match the visible declaration. \\
Compiler ABI lowering & Arguments and results use the target convention. \\
Assembler & Instruction operands and relocation fields are well formed. \\
Static linker & Selected definitions, ranges, and layout satisfy object contracts. \\
Loader & Segments and runtime dependencies can be established. \\
\bottomrule
\end{tabular}
\end{table}

\section{Retain a resolution explanation}
\label{sec:ch57-diagnostics}

For an undefined name, report at least the referencing object and spelling; when possible include relocation and source information.
For duplicate strong definitions, report both defining objects rather than only the second one encountered.
For archive extraction, a trace can state which unresolved name selected which member and what new demands that member introduced.
Such an explanation turns command-line-order debugging into a graph problem the reader can inspect.

Figure~\ref{fig:ch57-resolution} depicts a weak choice being replaced by a strong definition while local references remain outside this competition.
The winning identity must then be used consistently by every affected relocation.
Reusing a stale address cached when the weak definition was first seen would create a mixed program in which some references bind strongly and others weakly.
A clean phase boundary resolves identities first and computes their final addresses after layout.

\begin{figure}[htbp]
\centering
\begin{tikzpicture}[node distance=15mm and 22mm]
\node[basic block] (refs) {references to answer};
\node[basic block,below left=of refs] (w) {weak default};
\node[basic block,below right=of refs] (s) {strong definition};
\node[basic block,below=of s] (v) {one selected identity};
\draw[dashed,->] (refs)--node[left,font=\small,fill=white,inner sep=2pt] {candidate}(w);
\draw[cfg edge] (refs)--(s); \draw[cfg edge] (s)--(v);
\draw[back edge] (w)--node[below,font=\small] {superseded}(s);
\end{tikzpicture}
\caption{Binding selects one owning definition before layout assigns its address to all references.}
\label{fig:ch57-resolution}
\end{figure}

\section{Implementation boundaries and scaling}
\label{sec:ch57-implementation}

The working linker consumes explicitly named object files.
It implements the local/global/weak subset above, but not archive parsing or automatic member extraction.
The archive algorithm is exercised separately as a small model accompanying the chapter, so the distinction between implemented executable linking and a modeled extension remains visible.
Adding archives requires validating the container, its member sizes/alignment, and its name/index conventions before invoking the existing object parser.

Production linkers also need policies for duplicate sections, COMDAT groups, symbol versioning, and target-specific special symbols.
Their large symbol tables benefit from compact interned names, cached resolution state, and concurrent parsing.
Parallelism must preserve deterministic conflict diagnostics and output order where promised.
A faster implementation of an unspecified winner policy is still an unspecified linker.

\section{Exercises}
\label{sec:ch57-exercises}
\begin{enumerate}
\item\label{ex:ch57-local} Two objects each define local loop at offset zero and global helper in only the second. Describe the lookup key for all three identities.
\item\label{ex:ch57-weak} Process weak f, strong f, weak f in that order. Repeat with the strong definition first. State the selected owner and whether either order is an error.
\item\label{ex:ch57-group} Trace main, archive B, archive A for the worked archive example without a group, then with repeated group scanning.
\item\label{ex:ch57-extraction} A weak definition already satisfies f before an archive containing a strong f is visited. Explain why the strong definition need not replace it.
\item\label{ex:ch57-complexity} A group has 12 members and a full-summary sweep costs 100 operations. Give the coarse maximum sweep cost under Algorithm~\ref{alg:ch57-archive}, excluding object incorporation.
\end{enumerate}

\section{Further reading}
\label{sec:ch57-reading}
The ELF symbol chapter specifies binding, visibility, undefined entries, and common symbols~\cite{xinuos-2026-elf}.
GNU linker documentation specifies archive search commands and script behavior for that implementation~\cite{gnu-2026-ld}.
Levine's book is a broader reference on symbol management and libraries; its author's site also records errata~\cite{levine-1999-linkers}.
